\section{Formulas}

  % \item Number of permutations of length $n$ that have no fixed
  %     points (derangements): $D_0 = 1, D_1 = 0, D_n = (n - 1)(D_{n-1}
  %     + D_{n-2})$
  % \item Number of permutations of length $n$ that have exactly $k$
  %     fixed points: $\binom{n}{k} D_{n-k}$


  \begin{itemize}[leftmargin=*]
    \item \textbf{Legendre symbol:} $\left(\frac{a}{b}\right) = a^{(b-1)/2} \pmod{b}$, $b$ odd prime.
    \item \textbf{Heron's formula:} A triangle with side lengths
      $a,b,c$ has area $\sqrt{s(s-a)(s-b)(s-c)}$ where $s =
      \frac{a+b+c}{2}$.
    \item \textbf{Pick's theorem:} A polygon on an integer grid
      strictly containing $i$ lattice points and having $b$ lattice
      points on the boundary has area $i + \frac{b}{2} - 1$. (Nothing
      similar in higher dimensions)
    \item \textbf{Euler's totient:} The number of integers less than
      $n$ that are coprime to $n$ are $n\prod_{p|n}\left(1 - \frac{1}{p}\right)$
      where each $p$ is a distinct prime factor of $n$.
    \item \textbf{König's theorem:} In any bipartite graph $G=(L\cup R,E)$, the number
      of edges in a maximum matching is equal to the number of
      vertices in a minimum vertex cover. Let $U$ be the set of
      unmatched vertices in $L$, and $Z$ be the set of vertices that
      are either in $U$ or are connected to $U$ by an alternating
      path. Then $K=(L\setminus Z)\cup(R\cap Z)$ is the minimum
      vertex cover.
    \item A minumum Steiner tree for $n$ vertices requires at most $n-2$ additional Steiner vertices.
    \item The number of vertices of a graph is equal to its minimum
      vertex cover number plus the size of a maximum independent set.
    \item \textbf{Lagrange polynomial} through points $(x_0,y_0),\ldots,(x_k,y_k)$ is $L(x) = \sum_{j=0}^k y_j \prod_{\shortstack{$\scriptscriptstyle 0\leq m \leq k$ \\ $\scriptscriptstyle m\neq j$}} \frac{x-x_m}{x_j - x_m}$
    \item \textbf{Hook length formula:} If $\lambda$ is a Young diagram and $h_{\lambda}(i,j)$ is the hook-length of cell $(i,j)$, then then the number of Young tableux $d_{\lambda} = n!/\prod h_{\lambda}(i,j)$.
    \item \textbf{Möbius inversion formula:} If $f(n) = \sum_{d|n} g(d)$, then $g(n) = \sum_{d|n} \mu(d) f(n/d)$. If $f(n) = \sum_{m=1}^n g(\lfloor n/m\rfloor)$, then $g(n) = \sum_{m=1}^n \mu(m)f(\lfloor\frac{n}{m}\rfloor)$.
    \item \#primitive pythagorean triples with hypotenuse $<n$ approx $n/(2\pi)$.
    \item \textbf{Frobenius Number:} largest number which can't be
      expressed as a linear combination of numbers $a_1,\ldots,a_n$
      with non-negative coefficients. $g(a_1,a_2) = a_1a_2-a_1-a_2$,
      $N(a_1,a_2)=(a_1-1)(a_2-1)/2$. $g(d\cdot a_1,d\cdot a_2,a_3) =
      d\cdot g(a_1,a_2,a_3) + a_3(d-1)$. An integer $x>\left(\max_i
      a_i\right)^2$ can be expressed in such a way iff.\ $x\ |\
      \mathrm{gcd}(a_1,\ldots,a_n)$.
  \end{itemize}

  \subsection{Physics}
    \begin{itemize}
      \item \textbf{Snell's law:} $\frac{\sin\theta_1}{v_1} = \frac{\sin\theta_2}{v_2}$
    \end{itemize}

  \subsection{Markov Chains}
    A Markov Chain can be represented as a weighted directed graph of
    states, where the weight of an edge represents the probability of
    transitioning over that edge in one timestep. Let $P^{(m)} = (p^{(m)}_{ij})$
    be the probability matrix of transitioning from state $i$ to state $j$
    in $m$ timesteps, and note that $P^{(1)}$ is the adjacency matrix of
    the graph. \textbf{Chapman-Kolmogorov:} $p^{(m+n)}_{ij} = \sum_{k}
    p^{(m)}_{ik} p^{(n)}_{kj}$. It follows that $P^{(m+n)} =
    P^{(m)}P^{(n)}$ and $P^{(m)} = P^m$. If $p^{(0)}$ is the initial
    probability distribution (a vector), then $p^{(0)}P^{(m)}$ is the
    probability distribution after $m$ timesteps.

    The return times of a state $i$ is $R_i = \{m\ |\ p^{(m)}_{ii} > 0 \}$,
    and $i$ is \textit{aperiodic} if $\gcd(R_i) = 1$. A MC is aperiodic if
    any of its vertices is aperiodic. A MC is \textit{irreducible} if the
    corresponding graph is strongly connected.

    A distribution $\pi$ is stationary if $\pi P = \pi$. If MC is
    irreducible then $\pi_i = 1/\mathbb{E}[T_i]$, where $T_i$ is the
    expected time between two visits at $i$. $\pi_j/\pi_i$ is the expected
    number of visits at $j$ in between two consecutive visits at $i$. A MC
    is \textit{ergodic} if $\lim_{m\to\infty} p^{(0)} P^{m} = \pi$. A MC is
    ergodic iff.\ it is irreducible and aperiodic.

    A MC for a random walk in an undirected weighted graph (unweighted
    graph can be made weighted by adding $1$-weights) has $p_{uv} =
    w_{uv}/\sum_{x} w_{ux}$. If the graph is connected, then $\pi_u =
    \sum_{x} w_{ux} / \sum_{v}\sum_{x} w_{vx}$. Such a random walk is
    aperiodic iff.\ the graph is not bipartite.

    An \textit{absorbing} MC is of the form $P = \left(\begin{matrix} Q & R
    \\ 0 & I_r \end{matrix}\right)$. Let $N = \sum_{m=0}^\infty Q^m = (I_t
    - Q)^{-1}$. Then, if starting in state $i$, the expected number of
    steps till absorption is the $i$-th entry in $N1$. If starting in state
    $i$, the probability of being absorbed in state $j$ is the $(i,j)$-th
    entry of $NR$.

    Many problems on MC can be formulated in terms of a system of
    recurrence relations, and then solved using Gaussian elimination.

  \subsection{Burnside's Lemma}
    Let $G$ be a finite group that acts on a set $X$. For each $g$ in $G$
    let $X^g$ denote the set of elements in $X$ that are fixed by $g$. Then
    the number of orbits \[ |X/G| = \frac{1}{|G|} \sum_{g\in G} |X^g| \]

    \[
        Z(S_n) = \frac{1}{n} \sum_{l=1}^n a_l Z(S_{n-l})
    \]

  \subsection{Bézout's identity}
    If $(x,y)$ is any solution to $ax+by=d$ (e.g.\ found by the Extended
    Euclidean Algorithm), then all solutions are given by \[
    \left(x+k\frac{b}{\gcd(a,b)}, y-k\frac{a}{\gcd(a,b)}\right) \]

  \subsection{Misc}
    \subsubsection{Determinants and PM}
      \begin{align*}
        det(A) &= \sum_{\sigma \in S_n}\text{sgn}(\sigma)\prod_{i = 1}^n a_{i,\sigma(i)}\\
        perm(A) &= \sum_{\sigma \in S_n} \prod_{i = 1}^n a_{i,\sigma(i)}\\
        pf(A) &= \frac{1}{2^nn!}\sum_{\sigma \in S_{2n}} \text{sgn}(\sigma)\prod_{i = 1}^n a_{\sigma(2i-1),\sigma(2i)}\\ &= \sum_{M \in \text{PM}(n)} \text{sgn}(M) \prod_{(i,j) \in M} a_{i,j}
      \end{align*}

    \subsubsection{BEST Theorem}
      Count directed Eulerian cycles. Number of OST given by
      Kirchoff's Theorem (remove r/c with root) $\#\textsc{OST}(G,r)
      \cdot \prod_{v}(d_v-1)!$

    \subsubsection{Primitive Roots}
      Only exists when $n$ is $2, 4, p^k, 2p^k$, where $p$ odd prime. Assume
      $n$ prime. Number of primitive roots $\phi(\phi(n))$
      Let $g$ be primitive root. All primitive roots are of the form $g^k$
      where $k,\phi(p)$ are coprime.\\ $k$-roots:
      $g^{i \cdot \phi(n) / k}$ for $0 \leq i < k$

    \subsubsection{Sum of primes} For any multiplicative $f$:
      \[
          S(n,p) = S(n, p-1) - f(p) \cdot (S(n/p,p-1) - S(p-1,p-1))
      \]

    \subsubsection{Floor}
      \begin{align*}
          &\left\lfloor \left\lfloor x/y \right\rfloor / z \right\rfloor = \left\lfloor x / (yz) \right\rfloor \\
          &x \% y = x - y \left\lfloor x / y \right\rfloor
      \end{align*}
    \subsubsection{Large Primes}
      100894373, 103893941, 999999937, 1000000007, 16208191877


\iffalse
  \clearpage
  \section*{Practice Contest Checklist}
    \begin{itemize}
      \item How many operations per second? Compare to local machine.
      \item What is the stack size?
      \item How to use printf/scanf with long long/long double?
      \item Are \texttt{\_{}\_{}int128} and \texttt{\_{}\_{}float128} available?
      \item Does MLE give RTE or MLE as a verdict? What about stack overflow?
      \item What is \texttt{RAND\_{}MAX}?
      \item How does the judge handle extra spaces (or missing newlines) in the output?
      \item Look at documentation for programming languages.
      \item Try different programming languages: C++, Java and Python.
      \item Try the submit script.
      \item Try local programs: i?python[23], factor.
      \item Try submitting with \texttt{assert(false)} and \texttt{assert(true)}.
      \item Return-value from \texttt{main}.
      \item Look for directory with sample test cases.
      \item Make sure printing works.

      \item Remove this page from the notebook.
    \end{itemize}
\fi